% !TeX root = ../../main.tex

Die Experimente bestimmen, mit welcher Konfiguration eine Retry Strategie ausgeführt wird.
Eine Konfiguration besteht aus den folgenden Parametern.
\begin{itemize}
	\item Retry Count
	\item Retry Delay
	\item Retry Strategie
	\item Concurrency
	\item Workload (nicht variabel)
	\item Iterations (nicht variabel)
\end{itemize}
Der Retry Count legt fest, wie häufig eine Transaktion wiederholt werden darf.
Retry Delay beschreibt, wie lange eine Strategie initial warten muss, bevor sie die Transaktion wiederholen darf.
Durch Retry Strategie legt ein Experiment fest, mit welchen Strategien es den Retry ausführen soll.
Der Concurrency Parameter legt fest, wie viele gleichzeitige Anfragen eine Datenbank bekommt, die alle den SQL-Code
ausführen wollen.
Workload ist fest definiert, damit man vergleichbare Ergebnisse bekommt und Iterations sorgt dafür, dass die
Ergebnisse eine gewisse Stabilität haben und keine großen Ausreißer sind.
Durch diese Parameter lassen sich die Strategien testen, ohne dass es zu viel wird.

Der Aufbau einer Retry Strategie besteht daraus, dass sie SQL-Code ausführen soll.
Daraufhin bekommt sie von der Datenbank einen Status zurück, entweder ist die Transaktion erfolgreich gewesen oder
nicht.
Im nächsten Schritt wird dies getestet.
Gleichzeitig wird noch getestet ob der Retry Count überschritten wird oder nciht.
Ist die Überprüfung positiv, dass ein weiterer Retry möglich ist, wird der Retry angewandt.
Es wird der Retry Count erhöht und so lange gewartet, wie Delay besagt.
Danach wird die Transaktion erneut ausgeführt.

Dies geschieht so lange bis eine der Abbruchbedingungen zutrifft.
Entweder war die Transaktion erfolgreich, oder die Transaktion hat die maximal zulässige Anzahl an Wiederholungen
überschritten.
Ist das passiert, wird das Ergebnis zurückgegeben und in die Logdatei geschrieben.